\phantomsection
\chapter{Conclusiones}

    % Pendiente: conclusiones del trabajo

\phantomsection
\section{Líneas futuras}
\label{sec:lineas_futuras}

    Quedan fuera del alcance del PFI, y se plantean como líneas futuras, los siguientes puntos.

    \begin{itemize}
        \item La estratificación del riesgo, ya sea con los índices del sistema ABCD de Belin o con la predicción del riesgo de progresión mediante redes neuronales recurrentes. Ambas necesitan datos que CornOrb no incluye, como los valores ABCD o estudios de un mismo paciente en distintas visitas.

        \item La validación clínica con muestras locales de distintos centros de salud de la región, a partir de las solicitudes ya iniciadas (ver capítulo Desarrollo).

        \item La adaptación de modelos fundacionales al análisis de la topografía corneal. Requiere conjuntos de datos del segmento anterior del ojo y una capacidad de cómputo que exceden el PFI.

        \item La compatibilidad con otros equipos de diagnóstico, como Pentacam\textsuperscript{\scriptsize\textregistered}{} o Sirius\textsuperscript{\scriptsize\textregistered}{}, que depende de obtener sus archivos de exportación.

        \item La inscripción de KERALYTIX como producto médico ante la ANMAT y su lanzamiento comercial, cuya planificación se desarrolla en el capítulo Estudio financiero.
    \end{itemize}
